% =============================================================================
%  Resumen Capítulo 2 — PRE-PROCESAMIENTO
%  Libro: Evsukoff, A. (2020). Inteligência Computacional: Fundamentos e
%         Aplicações. Río de Janeiro: E-papers.
% =============================================================================
% =============================================================================
%  preambulo.tex — Preámbulo de los resúmenes en español
%  Evsukoff, A. (2020). Inteligência Computacional. Rio de Janeiro: E-papers.
% =============================================================================
\documentclass[11pt,a4paper]{article}

\usepackage[T1]{fontenc}
\usepackage[utf8]{inputenc}
\usepackage[spanish,es-tabla]{babel}
\usepackage{lmodern}
\usepackage[a4paper,margin=2.4cm]{geometry}
\usepackage{amsmath,amssymb,amsthm,mathtools}
\usepackage{bm}
\usepackage{enumitem}
\usepackage{xcolor}
\usepackage{tcolorbox}
\tcbuselibrary{skins,breakable}
\usepackage{hyperref}
\usepackage{microtype}
\usepackage{titlesec}
\usepackage{fancyhdr}
\usepackage{booktabs}
\usepackage{array}

\hypersetup{
  colorlinks=true,
  linkcolor=black!70,
  urlcolor=blue!60!black,
  citecolor=black!70,
  pdfborder={0 0 0}
}

\titleformat{\section}{\Large\bfseries\sffamily\color{black!85}}{\thesection}{0.7em}{}
\titleformat{\subsection}{\large\bfseries\sffamily\color{black!75}}{\thesubsection}{0.6em}{}

\pagestyle{fancy}
\fancyhf{}
\fancyhead[L]{\small\itshape Resumen — Inteligencia Computacional (Evsukoff, 2020)}
\fancyhead[R]{\small\thepage}
\renewcommand{\headrulewidth}{0.3pt}

\newtcolorbox{conceito}[1][]{
  colback=blue!4, colframe=blue!55!black, boxrule=0.6pt,
  arc=2mm, left=6pt, right=6pt, top=4pt, bottom=4pt,
  title={\bfseries Concepto clave}, fonttitle=\sffamily\bfseries,
  breakable, #1
}
\newtcolorbox{formula}[1][]{
  colback=gray!6, colframe=gray!50, boxrule=0.4pt,
  arc=1.5mm, left=6pt, right=6pt, top=3pt, bottom=3pt,
  breakable, #1
}
\newtcolorbox{aplicacao}[1][]{
  colback=green!4, colframe=green!45!black, boxrule=0.5pt,
  arc=2mm, left=6pt, right=6pt, top=3pt, bottom=3pt,
  title={\bfseries Aplicación}, fonttitle=\sffamily\bfseries,
  breakable, #1
}

\newcommand{\R}{\mathbb{R}}
\newcommand{\E}{\mathbb{E}}
\newcommand{\Var}{\operatorname{Var}}
\newcommand{\Cov}{\operatorname{Cov}}
\newcommand{\diag}{\operatorname{diag}}
\newcommand{\argmin}{\operatorname*{arg\,min}}
\newcommand{\argmax}{\operatorname*{arg\,max}}
\newcommand{\T}{^{\!\top}}
\DeclareMathOperator{\sign}{sign}
\DeclareMathOperator{\tr}{tr}


\title{\sffamily\bfseries Capítulo 2 — Pre-procesamiento\\
       \large Caracterización estadística, limpieza y transformaciones lineales}
\author{Resumen de estudio}
\date{}

\begin{document}
\maketitle
\thispagestyle{fancy}

\begin{conceito}
El pre-procesamiento abarca la \textbf{selección, limpieza y
transformación} de los datos, y suele consumir entre 60\% y 80\% del
tiempo de un proyecto de minería de datos. Su producto es el
\emph{conjunto de entrenamiento}: una matriz $\mathbf{X}\in\R^{N\times
p}$ (con eventual matriz de salida $\mathbf{Y}\in\R^{N\times q}$) cuyas
variables estén libres de \emph{outliers}, valores faltantes, escalas
incompatibles y redundancias graves.
\end{conceito}

\section{Conjuntos de datos y notación}

Un conjunto de datos estructurado es una tabla de $N$ registros y $p$
(o $p+q$) variables. Las variables pueden ser \emph{continuas},
\emph{discretas ordenadas} o \emph{categóricas (nominales)}; estas
últimas se codifican generalmente como variables indicadoras
(\emph{dummy}) binarias para que puedan ser procesadas por algoritmos
numéricos. El texto adopta la notación:
\[
\mathcal{T} = \{(\mathbf{x}(t),\mathbf{y}(t)),\; t=1,\dots,N\},\quad
\mathbf{x}(t)\in\R^p,\;\mathbf{y}(t)\in\R^q,
\]
con vectores en minúsculas-negrita, matrices en mayúsculas-negrita,
escalares en cursiva, conjuntos en letras caligráficas. Los problemas
pueden ser \emph{estáticos} (el orden de los registros es irrelevante),
\emph{dinámicos} (series temporales), \emph{georreferenciados} o
\emph{espacio-temporales}. El capítulo se concentra en el caso
estático.

\section{Caracterización estadística}

\subsection{Análisis univariable}

Para cada variable $x_i$ se estiman medidas de centralidad y
dispersión:
\begin{align*}
\bar{x}_i &= \frac{1}{N}\sum_{t=1}^{N} x_i(t), &
\hat{\sigma}_i^2 &= \frac{1}{N-1}\sum_{t=1}^{N}(x_i(t)-\bar{x}_i)^2,
\end{align*}
además de mínimo, máximo, mediana, percentiles $P_{25}$ y $P_{75}$
(cuartiles $Q_1,Q_3$), rango intercuartílico
($\mathrm{IRQ}=Q_3-Q_1$) y moda. La \textbf{media} puede ser engañosa
en distribuciones asimétricas o con \emph{outliers}; la \textbf{mediana}
es robusta. La discrepancia entre media y mediana ya sugiere asimetría.

La distribución normal $N(\mu,\sigma^2)$, con densidad
\[
f(x)=\frac{1}{\sqrt{2\pi\sigma^2}}\exp\!\left(-\frac{(x-\mu)^2}{2\sigma^2}\right),
\]
es fundamental: por el \emph{teorema central del límite}, las medias
muestrales de procesos estacionarios convergen en distribución a una
normal, independientemente de la distribución original. La función
gaussiana reaparece como \emph{función de base} en modelos RBF y SVM
(Capítulo~7).

El análisis visual emplea \textbf{histogramas} (incluyendo histogramas
por clase, en problemas de clasificación), \textbf{box-plots} y el
\textbf{gráfico de probabilidad normal} (la linealidad indica
normalidad). Histogramas multimodales pueden reflejar mezcla de clases
o subpoblaciones.

\subsection{Análisis multivariado}

La estadística multivariada-clave es la \textbf{matriz de covarianzas}
\begin{formula}
\[
\hat{\boldsymbol{\Sigma}} = \frac{1}{N-1}\sum_{t=1}^{N}
(\mathbf{x}(t)-\bar{\mathbf{x}})\T(\mathbf{x}(t)-\bar{\mathbf{x}})
\in\R^{p\times p},
\]
simétrica y definida positiva, con $\hat{\sigma}_{ii}^2 =
\hat{\sigma}_i^2$ en la diagonal.
\end{formula}

La versión estandarizada es la \textbf{matriz de correlación} de
Pearson, con elementos
\[
\hat{\rho}_{ij} = \frac{\Cov(x_i,x_j)}{\sqrt{\Var(x_i)\Var(x_j)}}\in[-1,1].
\]
Valores absolutos próximos a~1 indican dependencia lineal
(redundancia). Importante: $\hat{\rho}_{ij}=0$ excluye apenas
dependencia \emph{lineal}, no relación no lineal. La visualización
típica es la \emph{scatter plot matrix}, o, para muchas variables, un
mapa de color de la matriz de correlación.

\subsection{Distancias y similitudes}

Para detectar \emph{outliers}, imputar valores faltantes o hacer
\emph{clustering}, se calcula la matriz $\mathbf{D}\in\R^{N\times N}$
con $d_{ij}=\mathrm{dist}(\mathbf{x}(i),\mathbf{x}(j))$. Las
principales métricas (\emph{cualesquiera} que satisfagan no
negatividad, simetría y desigualdad triangular) son:

\begin{center}\small
\begin{tabular}{@{}l l@{}}
\toprule
Euclidiana ($L_2$) & $\sqrt{\sum_i (\nu_i-\upsilon_i)^2}$ \\
Manhattan ($L_1$) & $\sum_i |\nu_i-\upsilon_i|$ \\
Chebyshev ($L_\infty$) & $\max_i |\nu_i-\upsilon_i|$ \\
Minkowski ($L_r$) & $(\sum_i |\nu_i-\upsilon_i|^r)^{1/r}$ \\
Mahalanobis & $\sqrt{(\boldsymbol\nu-\boldsymbol\upsilon)\,
                     \hat{\boldsymbol{\Sigma}}^{-1}
                     (\boldsymbol\nu-\boldsymbol\upsilon)\T}$ \\
\bottomrule
\end{tabular}
\end{center}

Similitudes clásicas: coseno
$\sin_C(\boldsymbol\nu,\boldsymbol\upsilon)=
\langle\boldsymbol\nu,\boldsymbol\upsilon\rangle/
(\|\boldsymbol\nu\|\,\|\boldsymbol\upsilon\|)$ y gaussiana
$\sin_G=\exp(-\|\boldsymbol\nu-\boldsymbol\upsilon\|^2/(2\sigma^2))$. La
visualización de la matriz de similitudes como \emph{data image} (mapa
de calor) revela grupos y \emph{outliers} incluso en alta dimensión.

\section{Limpieza de datos}

\subsection{Estandarización}

\paragraph{Min--max:} $\hat{x}_i(t)=(x_i(t)-\min)/(\max-\min)\in[0,1]$.
Sensible a \emph{outliers}; sustituir por $P_5$ y $P_{95}$ si es
necesario.

\paragraph{$z$-score:} $\hat{x}_i(t)=(x_i(t)-\bar{x}_i)/\hat{\sigma}_i$.
Resulta en media nula y varianza unitaria; más robusto. En
distribuciones muy asimétricas conviene aplicar antes la transformación
logarítmica $\log(x+c)$.

\subsection{Detección de outliers}

Dos familias de métodos:
\begin{itemize}[leftmargin=1.2cm]
  \item \textbf{Por distribución:} bajo hipótesis de normalidad,
        rechazar registros fuera de $[\mu-3\sigma,\mu+3\sigma]$.
        Visualmente: \emph{box-plot}, con regla
        $[Q_1-1{,}5\,\mathrm{IRQ},\;Q_3+1{,}5\,\mathrm{IRQ}]$. Actúa
        \emph{variable a variable}; sólo señala sospechas.
  \item \textbf{Por distancia:} se ordenan las distancias medias
        $\bar{d}_i=\frac{1}{N}\sum_j d_{ij}$ y se descartan los
        $P_{out}\%$ registros más alejados (Algoritmo~2.1). Considera
        todas las variables simultáneamente; costo $\mathcal{O}(N^2)$.
\end{itemize}

\begin{aplicacao}
En detección de fraudes o anomalías, los \emph{outliers} \textbf{son}
el objeto de interés y no deben removerse.
\end{aplicacao}

\subsection{Valores faltantes}

Deben distinguirse del valor cero. Estrategias:
\begin{enumerate}[leftmargin=1.2cm]
  \item \textbf{Eliminar} registros (o variables) con valores
        faltantes — viable cuando la pérdida es pequeña o se concentra
        en pocas columnas;
  \item \textbf{Imputar}: sustituir por la media (ingenuo, pierde
        varianza) o, mejor, por los $k$-vecinos más próximos
        (Algoritmo~2.2): para cada registro incompleto, se busca el
        subespacio de las variables disponibles, se calculan distancias
        a los registros completos y se imputa la media (eventualmente
        ponderada por el inverso de la distancia) de los $k\geq 3$
        vecinos.
\end{enumerate}

La imputación introduce información artificial; debe medirse el impacto
sobre las estadísticas de las variables afectadas.

\section{Transformaciones lineales}

Una transformación lineal es una matriz $\mathbf{T}:\R^n\to\R^m$.
Cuando $\mathbf{T}$ es \emph{ortogonal}
($\mathbf{T}\T\mathbf{T}=\mathbf{I}$), la transformación preserva normas
y ángulos: se trata de una rotación del espacio.

\subsection{Análisis de Componentes Principales (ACP)}

\begin{conceito}
Encontrar una rotación ortogonal $\mathbf{P}$ tal que las variables
transformadas $\mathbf{Z}=\mathbf{X}\mathbf{P}$ sean
\textbf{decorrelacionadas} y ordenadas por varianza decreciente. Las
columnas de $\mathbf{P}$ son las \emph{componentes principales} (o
autovectores de la matriz de covarianzas).
\end{conceito}

La diagonalización de la matriz de covarianzas da origen al problema de
autovalores
\[
\hat{\boldsymbol{\Sigma}}\,\mathbf{p}_i = \lambda_i\,\mathbf{p}_i,
\quad
\hat{\boldsymbol{\Sigma}}\,\mathbf{P} = \mathbf{P}\,\boldsymbol\Lambda,
\]
con $\boldsymbol\Lambda=\diag(\lambda_1,\dots,\lambda_p)$ y
$\lambda_1\geq\dots\geq\lambda_p\geq 0$. El \textbf{cociente de
Rayleigh}
\[
\lambda_{\max}=\max_{\mathbf{v}\neq 0}\frac{\mathbf{v}\T\hat{\boldsymbol{\Sigma}}\mathbf{v}}
                                            {\mathbf{v}\T\mathbf{v}}
\]
muestra que la primera componente es la \textbf{dirección de máxima
varianza}; las demás son las direcciones ortogonales subsecuentes de
máxima varianza.

La invariancia de la traza bajo rotaciones ortogonales garantiza
$\sum_i \hat{\sigma}_i^2 = \sum_i \lambda_i$, de modo que la fracción
de varianza explicada por las $n$ primeras componentes es
$\sum_{i=1}^n \lambda_i/\sum_{i=1}^p \lambda_i$. Típicamente se elige
$n$ tal que esa fracción sea $\geq 85\%$. Se reduce así la
dimensionalidad preservando lo esencial de la estructura lineal.

\paragraph{Pre-requisito:} centrar/estandarizar las variables
($z$-scores) antes del cálculo, pues la ACP es sensible a la escala.
También supone que los datos sean aproximadamente gaussianos.

\subsection{Descomposición en Valores Singulares (SVD)}

Para cualquier matriz $\mathbf{A}\in\R^{n\times m}$ existe la
factorización
\[
\mathbf{A}=\mathbf{U}\,\mathbf{S}\,\mathbf{V}\T,
\]
con $\mathbf{U},\mathbf{V}$ ortogonales y $\mathbf{S}$ diagonal con
valores singulares $s_1\geq\dots\geq s_r\geq 0$. Aplicada a la matriz
de datos $\mathbf{X}$, se obtiene
$\mathbf{X}\T\mathbf{X}=\mathbf{V}\mathbf{S}^2\mathbf{V}\T$, es decir,
$\mathbf{V}=\mathbf{P}$ y $s_i=\sqrt{\lambda_i}$. La SVD permite
calcular la ACP \textbf{sin} formar el producto
$\mathbf{X}\T\mathbf{X}$, ganando estabilidad numérica, y tiene
aplicaciones propias en \emph{Latent Semantic Indexing} (LSI) y
\emph{recuperación de información}.

\subsection{Análisis Discriminante Lineal (ADL)}

La ACP ignora la variable de salida; en clasificación, la mejor
proyección puede \emph{no} ser la de mayor varianza. El ADL (Fisher,
1936) busca la transformación $\mathbf{V}$ que maximiza la separación
entre clases minimizando la dispersión intra-clases, vía descomposición
de la varianza total
\[
\boldsymbol\Sigma_T = \boldsymbol\Sigma_B + \boldsymbol\Sigma_W,
\]
en varianza \emph{between} e \emph{within} clases. Se resuelve el
problema de autovalor generalizado
\[
\boldsymbol\Sigma_B\,\mathbf{V} = \boldsymbol\Sigma_W\,\mathbf{V}\,
\boldsymbol\Lambda,
\]
generando como máximo $m-1$ direcciones discriminantes (con
$m$~clases). Los datos proyectados son $\mathbf{Z}=\mathbf{X}\mathbf{V}$.

\subsection{Análisis de Correlación Canónica (ACC)}

Aborda problemas de regresión con dos conjuntos de variables,
$\mathbf{X}\in\R^{N\times p}$ e $\mathbf{Y}\in\R^{N\times q}$. Calcula
simultáneamente dos transformaciones $\mathbf{U}_x,\mathbf{U}_y$ que
maximizan la correlación cruzada entre las variables canónicas
$\mathbf{Z}_x=\mathbf{X}\mathbf{U}_x$ y
$\mathbf{Z}_y=\mathbf{Y}\mathbf{U}_y$. Algorítmicamente, se factorizan
$\mathbf{X}=\mathbf{Q}_x\mathbf{R}_x$ e
$\mathbf{Y}=\mathbf{Q}_y\mathbf{R}_y$ (descomposición QR) y se calcula
la SVD de $\mathbf{Q}_x\T\mathbf{Q}_y$; las correlaciones canónicas son
los valores singulares.

\section{Aplicaciones ilustrativas}

\begin{aplicacao}
\textbf{Patrones de movilidad urbana:} a partir de 2{,}1 mil millones de
registros de llamadas (\emph{Call Detail Records}, CDR) en la Región
Metropolitana de Río de Janeiro (2014), fue posible estimar matrices
origen--destino e identificar dos patrones claros — viajes
\emph{casa-trabajo} (Centro--Zona Sur, picos en días hábiles, simétrico
ida/vuelta) y viajes \emph{turísticos} (Zona Sur--Teresópolis, picos en
fines de semana y feriados, desfase ida/vuelta).
\end{aplicacao}

\begin{aplicacao}
\textbf{Validación de sensores por ecuaciones de paridad:} en un
sistema lineal estático con $p$ sensores y $n$ estados ($p>n$), las
columnas de la matriz de componentes principales se descomponen como
$\mathbf{P}=[\mathbf{C}\mid\mathbf{W}]$. El vector de paridades
$\mathbf{u}(t)=\mathbf{x}(t)\mathbf{W}$, satisfaciendo
$\mathbf{C}\T\mathbf{W}=\mathbf{0}$, se anula en ausencia de fallas e
indica direcciones típicas de falla cuando es perturbado. Un umbral
$\xi=\sqrt{\sum_{i=n+1}^{p}\lambda_i}$ acomoda el ruido. Aplicado a la
hidroeléctrica de Funil, detectó fallas en extensómetros incluso con
muchos valores faltantes.
\end{aplicacao}

\section{Comentarios finales}

\begin{enumerate}[leftmargin=1.2cm]
  \item Invertir tiempo en la \textbf{caracterización y visualización}
        de los datos es la mejor garantía de un modelado exitoso en las
        etapas siguientes.
  \item Las \textbf{transformaciones lineales} (ACP, SVD, ADL, ACC) son
        herramientas universales para reducir dimensionalidad,
        decorrelacionar variables o alinear la representación con la
        tarea supervisada. Todas dependen fuertemente de una
        \textbf{estandarización} previa.
  \item Los métodos basados en \textbf{distancias y similitudes} son lo
        suficientemente generales para resolver otros problemas
        (limpieza, \emph{outliers}, imputación) que aparecen
        repetidamente en el libro.
  \item Existen extensiones no lineales de los mismos principios (ICA,
        LLE, ACP por funciones de núcleo); el capítulo se limita al
        caso lineal, que es la base obligatoria para los modelos no
        lineales de los Capítulos~7--9.
\end{enumerate}

\end{document}
